\title{The \textsf{Y} Language\\[2pt]\large A Formal Definition}
\author{the y compiler project}
\date{October 2026}

\maketitle

\begin{abstract}
This document is the language definition of \textsf{Y}, a small
statically-checked systems language that compiles directly to static
x86-64 executables --- Linux ELF by default, a Windows PE, a freestanding
kernel image, or a WebAssembly module for the web --- on request. It records the complete surface language, its type
system, and its evaluation semantics --- with particular precision for the
mutating array operators (\ic{+=}, \ic{*=}), the
increment/decrement and compound-assignment families, and their
evaluated-exactly-once rules, so that these features are specified rather
than discovered. Where the language is quirky, the quirk is written
down; where a construct is reserved but unimplemented, it says so.
\end{abstract}

\tableofcontents

\section{Overview}
\label{sec:overview}

\textsf{Y} is a small imperative language with a static checker, designed
so that a complete implementation of the language fits in one head: the
compiler is written in \textsf{Y} itself (\ic{src/y/}) and compiles a
program to a \emph{static} x86-64 ELF executable using only
its own assembler and ELF writer --- no system assembler, linker, or C
toolchain is involved at run time or at build time of user programs.

Two further targets come from the same pipeline: \ic{--target=win}
writes a PE32+ executable (Section~\ref{sec:toolchain}) and
\ic{--target=kernel} writes a freestanding image that boots on bare metal
(Section~\ref{sec:os}).

\subsection*{Implementations}

The compiler in \ic{src/y/} is self-hosting and byte-reproducible.
A generated C seed (\ic{bootstrap/ycc.c}) lets GCC build it without a Y
executable; \ic{src/c/} contains the seed's boxed value runtime.
Compiler passes are maintained in Y.

This document defines observable behavior, including value sharing,
integer wraparound and trap conditions.

\subsection*{Compiler options}

\ic{ycc [OPTIONS] FILE}, \ic{ycc run [OPTIONS] FILE [-- ARG...]}, and
\ic{ycc build [OPTIONS] [DIR]} share their code-generation options.
\ic{--optimize=speed} (the default) and \ic{--optimize=size} select mutually
exclusive goals; there are no numeric optimization levels. Optimizations that
improve both code size and execution speed run in both modes. These include
removing redundant moves, branches to the following label, and dead instructions,
and using zero idioms where arithmetic flags are dead. Speed mode duplicates
short return sequences to remove taken branches; size mode uses compact
constant loads. These final-link passes also optimize bundled standard-library
units and cached packages. Neither mode changes language semantics.

\ic{--emit=asm|source|wat} selects assembly, expanded source, or a wasm
listing (the executable is the default). \ic{--output=PATH} (or \ic{-o PATH})
chooses an executable or assembly path; assembly and the wasm listing default
to stdout, and expanded source goes to stdout. \ic{--target=linux|win|kernel|web}
names the OS and \ic{--arch=x64|aarch64} the instruction set (Linux only; the
default arch is the host's). \ic{--debug} (or \ic{-g}) includes symbols and
source lines.
\ic{--obfuscate[=MODE]} is an opt-in code-protection mode; it never changes what a
program does, and there is one level: the bare flag and \ic{--obfuscate=full}
select the strongest preset. The program's own
bodies are rewritten at final link: numeric constants are rebuilt, stack frames
are shifted, dead bodies, opaque guards and junk bytes are added, calls may
become indirect, basic blocks are reordered, and the diversity-default static
pass is applied. The image is encrypted and starts through a plaintext page;
code pages stay encrypted until first use (a page fault decodes one page), and
eligible helpers, including float/SSE and division shapes, become generated
bytecode run by a per-build interpreter with a per-function opcode map.
Runtime messages are decoded eagerly because a system call reading a protected
page fails with EFAULT instead of taking the page-fault handler; literals and
the bytecode are decoded page by page, and a literal use touches one byte per
page first. \ic{--obfuscate-seed=N} makes a build reproducible;
\ic{--obfuscate-key=PATH} keeps the stream key in a 16-byte external file;
the sub-flags \ic{--obfuscate-vm=off|auto|all|bytecode} (default \ic{bytecode}),
\ic{--obfuscate-layout=source|shuffle}, \ic{--obfuscate-std=off|auto|all} and
\ic{--obfuscate-diversity=standard|high} tune a layer for tests and bring-up;
\ic{--obfuscate-no-anti-debug} disables the dumpable/tracer checks;
\ic{--obfuscate-report} prints a per-build id and counts;
\ic{--obfuscate-map=PATH} writes a label/address map outside the binary. The
cipher and loader are per-build and page-keyed, with
runtime-dependent opaque branches, jumps over random bytes and a key-derived
startup page schedule; there is no plaintext per-function table. A keyed hash
of the sealed ciphertext refuses a modified file, a keyed per-page MAC refuses
a page whose plaintext is wrong when it first decrypts, and a key-derived
canary replaces the old fixed magic. Full mode reads TracerPid instead of becoming a
tracee, so lazy page faults stay with the loader. Neither makes a standalone
binary cryptographically secure.
Obfuscation covers Linux ELF and Windows PE; a PE has a plaintext entry stub
that decrypts code and literals in place in a writable \ic{.text}, uses the
flatten layer and the embedded key, and has no lazy pages or per-page MACs.
It refuses \ic{-g}, \ic{--pack}, \ic{--emit=asm}, \ic{--emit=source} and \ic{--target=kernel}. The mode raises the cost of static and
dynamic analysis; it is not a cryptographic guarantee.
\ic{--warnings=none|default|extra|all|fix} selects diagnostics by name
(\ic{fix} applies the safe shortenings in place). \ic{--gc=never|auto|always|stress},
\ic{--regions=auto|off}, \ic{--exec=on|off}, and \ic{--std=pack|source} control
the collector, automatic regions, the \ic{exec()} builtin, and standard-unit
sources. \ic{--pack[=METHOD]} writes a self-decompressing executable and
\ic{--time} prints a phase table. Value options accept either
\ic{--name=value} or \ic{--name value}; \ic{-O}, \ic{-o} and \ic{-p} are the
short forms. Build chooses its own output paths and accepts executable output
only.

The default is \ic{--gc=never}: executables contain no collector. The
lifetime solver places temporary allocations in regions and rejects a
program when it cannot prove reclamation for repeated allocations.
\ic{--gc=auto} explicitly permits collection where needed;
\ic{--gc=always} forces it; \ic{--gc=stress} runs a cycle at every allocation.

\subsection*{Notation}

Grammar productions are given in EBNF, with
\texttt{typewriter} for terminals, \ty{italic} for nonterminals,
\texttt{[~]} for option, \texttt{\{~\}} for repetition,
\texttt{|} for alternation, and \texttt{(~)} for grouping.
Semantic rules are written as equational desugars or prose; the desugars
are normative: an implementation may use any mechanism that produces the
same observable behavior.

\subsection*{A taste}

\begin{lstlisting}
// values, references, and the mutating operators
let xs = [1, 2, 3];      // heap array: len 3
xs += 4;                 // in-place push:        [1, 2, 3, 4]
xs *= 0;                 // in-place prepend:     [0, 1, 2, 3, 4]
let ys = xs + [9];       // NEW array, xs untouched

// the target of a compound assignment runs exactly once
let a = [1, 2];
def idx() -> long { print("called"); 0 };
a[idx()] += 10;          // idx() runs ONE time

// records are reference values too
let r = { name: "y", nums: [1, 2] };
r.nums += 3;             // visible through every binding of r
\end{lstlisting}

\section{Lexical structure}
\label{sec:lexical}

\subsection{Source text}

A program is a sequence of bytes, conventionally UTF-8. The lexer works
on bytes; only ASCII is significant to the grammar. Both \ic{"\bs n"} (LF) and \ic{"\bs r\bs n"} (CRLF) end a line. Positions reported
in diagnostics are 1-based line and column, the column counted in bytes.

\subsection{Comments}

\begin{lstlisting}
// a line comment runs to the end of the line
/* a block comment
   does not nest: the first */ closes it */
\end{lstlisting}

\noindent Comments are whitespace: they separate tokens and nothing more.

\subsection{Identifiers and keywords}

\ic{isalpha}(c) is true for \ic{A}--\ic{Z}, \ic{a}--\ic{z} and
\ic{\_}; \ic{isdigit}(c) for \ic{0}--\ic{9}. An identifier is

\begin{align*}
\ty{identifier} &::= \ty{isalpha}~\{~\ty{isalnum}~\}
\end{align*}

\noindent so \ic{\_} is a word character: \ic{is\_digit} is a single
identifier, not three tokens. The prefix \ic{y\_} is \emph{reserved}:
identifiers beginning with \ic{y\_} are used by the compiler's generated
code and must not be written by hand.

The keywords are:
\begin{center}
\ic{let} \ic{const} \ic{def} \ic{if} \ic{else} \ic{while} \ic{for}
\ic{foreach} \ic{in} \ic{return} \ic{true} \ic{false}
\end{center}

\subsection{Literals}
\label{sec:literals}

\paragraph{Integers.} Decimal, binary (\ic{0b1010}) and hexadecimal
(\ic{0x1F}) radix literals are supported. Any of them may group digits
with \ic{\_}, which is ignored: a radix literal (\ic{0x} / \ic{0b},
either case) in any position after the prefix (\ic{0xFF\_FF},
\ic{0b1010\_1010}), a decimal one anywhere after its first digit
(\ic{1\_000\_000}; a leading \ic{\_} starts an identifier instead).
A prefix needs a digit of its radix: \ic{0x}, \ic{0b\_} and \ic{0b2}
are errors (\ic{binary literal '0b2' has no digits}), one for the whole
word.
There is no octal form: \ic{017} is decimal seventeen. All integer
literals are \ty{long} (64-bit). A radix literal spells the 64 bits
themselves, so a mask may set the top bit:
\ic{0xFFFF\_FFFF\_FFFF\_FFFF} is \ic{-1} and
\ic{0x8000\_0000\_0000\_0000} the smallest \ty{long}; digits past 64
bits wrap away (modulo $2^{64}$).

\paragraph{Floats.} A digit sequence with a decimal point and/or an
exponent: \ic{1.5}, \ic{1.5e+3}, \ic{2.0e-1}. \ic{\_} may group digits
after the first digit of the whole part, of the fraction and of the
exponent: \ic{1\_000.5}, \ic{0.000\_001}, \ic{1e1\_0} (\ic{1.\_5} and
\ic{1e\_5} are parse errors). A dot followed by a non-digit is
\emph{not} a fractional part --- \ic{a..b} lexes as \ic{a}, \ic{..},
\ic{b}. Float literals are IEEE-754 binary64 (\ty{float}).

\paragraph{Strings.} Double-quoted byte strings:

\begin{center}
\begin{tabular}{ll}
\toprule
escape & meaning \\
\midrule
\ic{\bs n} \ic{\bs t} \ic{\bs r} & LF, HT, CR \\
\ic{"\bs "} \ic{"\bs ""} \ic{"\bs '"} & backslash, quote \\
\ic{\bs xHH} & byte with hex value \ic{HH} \\
\ic{\bs c} (any other \ty{c}) & \ty{c}; the backslash is dropped \\
\bottomrule
\end{tabular}
\end{center}

\noindent An escape is read as a single byte, and an unrecognised
\ic{\bs c} simply drops the backslash (\ic{"q\bs q"} is \ic{"qq"}), so a
NUL is spelled \ic{"\bs x00"} --- which is not \ic{"\bs0"} (that is the
character \ic{0}). A \ic{str} is a length-prefixed byte string, so an
escaped NUL inside one is ordinary data and \ic{len} counts it (use
\ic{[byte]} arrays, Section~\ref{sec:builtin-bytes}, for bulk binary
data).

\paragraph{Byte strings.} \ic{b"..."} and \ic{b'...'} are \emph{byte
string literals}: a \ic{[byte]} value spelled like a Python \ic{bytes}
literal. The escapes are the string escapes above, so binary data writes
naturally: \ic{b"\bs x7fELF"} is the four-byte ELF magic. Either quote
spells the same type; single quotes exist only for byte strings.

\paragraph{Format strings.} \ic{f"..."} is a string literal with
\emph{fields}, as in Python:
\begin{center}
\ic{f"\{name\} is \{age\} (\{share:.1\%\})"}
\end{center}
\noindent A field \ic{\{e\}} becomes \ic{tostring(e)}, and
\ic{\{e:spec\}} becomes \ic{e} formatted as \ty{spec} says
(Section~\ref{sec:format-spec}). The text around the fields takes the
string escapes, and \ic{\{\{} and \ic{\}\}} (or \ic{\bs\{} and
\ic{\bs\}}) are literal braces. An f-string is sugar for concatenation:
\ic{f"a\{x\}b"} means \ic{"a" + tostring(x)} \ic{+ "b"}, a \ty{str} like
any other.

A field holds one expression --- calls, string literals, record literals
and nested f-strings included --- and may span lines; an error inside it
points at its own line and column. The expression ends at the first
\ic{\}} or \ic{:} outside brackets, except a \ic{:} that answers a
\ic{?}: \ic{\{ok ? a : b\}} works, while a cons has to be parenthesised,
\ic{\{(x :: xs)\}}. A \ic{?} that is part of a \ic{?.} (Section~
\ref{sec:opt-propagation}) answers nothing, so the \ic{:} after it is
the spec's own: \ic{f"\{o?.port:x\}"\} is \ic{o?.port} as \ic{x}. The spec
is plain text up to the field's \ic{\}}; it cannot hold another field.

\paragraph{Booleans and none.} \ic{true}, \ic{false}. The empty option
is spelled \ic{none}, a predeclared name rather than a keyword
(Section~\ref{sec:type-options}). There is no \ic{null}: the word is an
ordinary (undefined) identifier.

\subsection{Operators and punctuation}

\begin{center}
\begin{tabular}{lll}
\toprule
kind & tokens \\
\midrule
% `<{}<`: the {} keeps T1's << and >> ligatures (guillemets) out.
three- and four-character &
  \ic{>{}>{}>} \ic{<{}<=} \ic{>{}>=} \ic{>{}>{}>=} \\
two-character &
  \ic{==} \ic{!=} \ic{<=} \ic{>=} \ic{\&\&} \ic{||} \ic{::} \ic{..}
  \ic{->} \ic{=>} \ic{++} \ic{--} \ic{<{}<} \ic{>{}>} \\
 &
  \ic{+=} \ic{-=} \ic{*=} \ic{/=} \ic{\%=}
  \ic{\&=} \ic{|=} \ic{\textasciicircum{}=} \\
single-character &
  \ic{+} \ic{-} \ic{*} \ic{/} \ic{\%} \ic{<} \ic{>} \ic{=} \ic{!}
  \ic{\&} \ic{.} \ic{?} \ic{|} \ic{\textasciicircum{}} \ic{\textasciitilde{}} \\
punctuation &
  \ic{(} \ic{)} \ic{\{} \ic{\}} \ic{[} \ic{]} \ic{,} \ic{;} \ic{:} \\
\bottomrule
\end{tabular}
\end{center}

\paragraph{Maximal munch.} The lexer always takes the longest operator
that matches, exactly like C:
\begin{itemize}[nosep]
\item \ic{x+++y} is \ic{x ++ + y}, not \ic{x + ++y};
\item \ic{x--y} is \ic{x-- y} (a parse error), not \ic{x - -y};
      write \ic{x - -y} when negation is meant;
\item \ic{/=} is one token; \ic{//} starts a comment, and
      \ic{/ == y} must be written with a space to divide;
\item \ic{x>{}>{}>=1} is the compound \ic{x >{}>{}>= 1};
\end{itemize}

\paragraph{Contextual keywords.} \ic{none}, \ic{match}, \ic{type} and
\ic{void} are ordinary identifiers to the lexer: \ic{none} is a
predeclared name (Section~\ref{sec:type-options}), \ic{match} is the
match-expression keyword only when an expression and \ic{\{} follow it,
and \ic{type} starts a type alias only when followed by \ic{name =}. Each
stays usable as an ordinary name elsewhere.

\section{Syntax}
\label{sec:grammar}

\subsection{Structure and modules}

A \ty{program} is a sequence of \ty{statement}s executed in order.
Source files may be combined with imports; an import splices the
imported file's text in front of the importing file's own code, under a
namespace alias:

\begin{lstlisting}
import io from "./std/io.y";   // io.f(...) addresses defs of io.y
\end{lstlisting}

\noindent An \ic{export \{ a, b \}} clause, when present, must be the
\emph{last} statement of a file. Paths resolve relative to the importing
file's directory.

\subsection{Declarations}

\begin{align*}
\ty{declaration} ::={}& \ic{let}~\ty{name}~[\ic{:}~\ty{type}]~\ic{=}~\ty{expr}~\ic{;} \\
  |~& \ic{const}~\ty{name}~[\ic{:}~\ty{type}]~\ic{=}~\ty{expr}~\ic{;} \\
  |~& \ic{def}~\ty{name}~\ic{(}~\ty{params}~\ic{)}~[\ic{->}~\ty{type}]~\ty{block}
\end{align*}

\ty{params} is a comma-separated list of \ty{name}~\ic{:}~\ty{type};
parameter annotations are \textbf{mandatory}, return types are optional
(inferred from the body, Section~\ref{sec:types}). A \ic{const}
declaration currently behaves exactly like \ic{let}; the keyword records
intent and is reserved for future immutability checking.

Function literals may also appear bound or anonymous:
\ic{let f = def (x: long) -> long \{ x \};}.

\subsection{Types}
\label{sec:type-grammar}

\begin{align*}
\ty{type} ::={}& \ic{long} \mid \ic{byte} \mid \ic{bool} \mid \ic{float}
              \mid \ic{str} \\
  |~& \ic{[}~\ty{type}~\ic{]}                    && \text{(heap array)}\\
  |~& \ic{[}~\ty{type}~\ic{;}~\ty{nat}~\ic{]}   && \text{(fixed array)}\\
  |~& \ty{type}~\ic{?}                           && \text{(option)}\\
  |~& \ic{*}~\ty{type}                           && \text{(pointer)}\\
  |~& \ic{(}~[\ty{type}~\{,\ty{type}\}]~\ic{)}~\ic{->}~\ty{type}
                                                 && \text{(function)}\\
  |~& \ic{\{}~[\ty{name}~\ic{:}~\ty{type}
        ~\{,~\ty{name}~\ic{:}~\ty{type}\}]~\ic{\}}
                                                 && \text{(record shape)}
\end{align*}

\noindent A record shape is structural; the shape \ic{\{ k: str \}} is
written to mean ``a record with at least a field \ic{k} of type
\ty{str}'' and is the conventional shape-hole annotation for AST-like
data. \ic{nat} is a non-negative integer literal. The type name
\ic{int} is accepted as a synonym of \ic{long} in annotations.

\subsection{Statements}

\begin{align*}
\ty{statement} ::={}& \ty{declaration} \\
  |~& \ty{expr}~\ic{;} \\
  |~& \ty{assignment}~\ic{;} \\
  |~& \ic{return}~[\ty{expr}]~\ic{;} \\
  |~& \ty{block} \\
  |~& \ic{if}~\ic{(}~\ty{expr}~\ic{)}~\ty{block}~[\ic{else}~(\ty{block} \mid \ty{statement})] \\
  |~& \ic{while}~\ic{(}~\ty{expr}~\ic{)}~\ty{block} \\
  |~& \ic{for}~\ic{(}~\ty{stmt}~\ic{;}~\ty{expr}~\ic{;}~\ty{stmt}~\ic{)}~\ty{block} \\
  |~& \ic{for}~\ty{name}~\ic{in}~\ty{expr}~\ic{..}~\ty{expr}~\ty{block} \\
  |~& \ic{foreach}~\ic{(}~\ty{name}~\ic{in}~\ty{expr}~\ic{)}~\ty{block} \\
  |~& \ty{block}
\end{align*}

\noindent \ty{block} is \ic{\{}~\ty{statement}\(^*\)~[\ty{expr}]~\ic{\}}:
a block is an \emph{expression} whose value is its trailing expression
(if any). A block must never be empty; write \ic{\{ 0; \}} where a
placeholder is needed. \ic{foreach} iterates the elements of an array;
the \ic{for} range form iterates \ic{lo..hi} with the low bound
inclusive and the high bound exclusive.

\subsection{Expressions}
\label{sec:expr-grammar}

\begin{align*}
\ty{primary} ::={}& \ty{long} \mid \ty{float} \mid \ty{string}
     \mid \ic{true} \mid \ic{false} \mid \ic{null} \mid \ty{name} \\
  |~& \ic{(}~\ty{expr}~\ic{)} \mid \ty{block} \\
  |~& \ic{[}~[\ty{expr}~\{,\ty{expr}\}]~\ic{]}
       \mid \ic{[}~\ty{expr}~\ic{;}~\ty{expr}~\ic{]} \\
  |~& \ic{\{}~[\ty{name}~\ic{:}~\ty{expr}~\{,\ldots\}]~\ic{\}}
       \quad \ic{\{}~[\ty{string}~\ic{:}~\ty{expr}~\{,\ldots\}]~\ic{\}} \\
  |~& \ic{def}~\ic{(}~\ty{params}~\ic{)}~[\ic{->}~\ty{type}]~\ty{block}
\end{align*}

\begin{align*}
\ty{postfix} ::={}& \ty{primary}
   ~\{~\ic{(}~[\ty{args}]~\ic{)} \mid \ic{[}~\ty{expr}~\ic{]}
   \mid \ic{[}~\ty{expr}\,?\ic{:}\,\ty{expr}\,?~\ic{]} \mid \ic{.}~\ty{name}
   \mid \ic{++} \mid \ic{--}~\} \\
\ty{unary} ::={}& [\ic{!} \mid \ic{-} \mid \ic{*} \mid \ic{\&}
   \mid \ic{++} \mid \ic{--}]~\ty{postfix}
\end{align*}

\subsection*{Operator precedence}

All binary operators are left-associative; lower rows bind tighter:

\begin{center}
\begin{tabular}{cll}
\toprule
level & operators & notes \\
\midrule
1 & \ic{||} & short-circuit \\
2 & \ic{\&\&} & short-circuit \\
3 & \ic{==} \ic{!=} & structural equality \\
4 & \ic{<} \ic{>} \ic{<=} \ic{>=} & numeric, or lexicographic \ty{str} \\
5 & \ic{::} & one-element cons (Section~\ref{sec:arrays}) \\
6 & \ic{+} \ic{-} & numeric, string/array \ic{+} \\
7 & \ic{*} \ic{/} \ic{\%} & numeric \\
\bottomrule
\end{tabular}
\end{center}

\subsection*{Assignment forms}

\begin{align*}
\ty{assignment} ::={}& \ty{place}~\ic{=}~\ty{expr} \\
  |~& \ty{place}~\ty{compound-op}~\ty{expr}
     && \ic{+=}~\ic{-=}~\ic{*=}~\ic{/=}~\ic{\%=} \\
  |~& \ty{place}~\ic{++} \mid \ic{++}~\ty{place}
   \mid \ty{place}~\ic{--} \mid \ic{--}~\ty{place}
\end{align*}

\noindent A \ty{place} is a name, an index expression \ic{a[e]} (arrays
and maps), or a field expression \ic{o.f} (records), possibly nested:
\ic{grid[0][1]}, \ic{rec.v[2]}, \ic{m['k'][i]}. Field access is
symbolic and exclusive: \ic{r.f} for records, \ic{m['k']} for maps.
Pointers add \ic{*p = e}.
Section~\ref{sec:exactly-once} defines the exact evaluation of every
form; all of them evaluate the place's subexpressions \textbf{exactly
once}.

\section{Types and static semantics}
\label{sec:types}

Checking is static and \emph{fatal}: a type error is a compile error.
The type system is small --- scalars, arrays, records, maps, options,
pointers, functions and \ic{void} --- and mostly structural.

\subsection{Scalar types}

\begin{center}
\begin{tabular}{lll}
\toprule
type & domain & notes \\
\midrule
\ic{long}  & 64-bit signed integers & wraparound on overflow \\
\ic{byte}  & $0\ldots255$ & stored in one byte inside \ic{[byte]} \\
\ic{bool}  & \ic{true}, \ic{false} & \\
\ic{float} & IEEE-754 binary64 & \\
\ic{str}   & immutable byte strings & length-prefixed; may hold NUL \\
\bottomrule
\end{tabular}
\end{center}

\subsection{Arrays}
\label{sec:type-arrays}

\ic{[T]} is a heap array of \ty{T}, \ic{[T; N]} a fixed-length array of
\ty{T} with \ty{N} a literal. The element type is part of the type and
arrays are \textbf{homogeneous}: every element of an array has exactly
the element type. Heterogeneous array literals are rejected:

\begin{lstlisting}
let bad = [1, "two"];        // error: array element: expected long, got str
\end{lstlisting}

\noindent This applies recursively --- \ic{[[1], [2]]} has type
\ic{[[long]]}, and every operator that mixes arrays and elements
(Section~\ref{sec:arrays}) enforces the same rule.

\subsection{Records}
\label{sec:type-records}

A record literal \ic{\{ a: 1, b: "x" \}} has the structural shape
\ic{\{ a: long, b: str \}}; field order is preserved and is part of the
type. A field read is legal when the static shape has the field. The
shape-hole form \ic{\{ k: str \}} matches any record with at least those
fields and is used where the precise shape is not known statically (it
poisons quietly rather than erroring --- see
Section~\ref{sec:poison}).

\subsection{Maps}

\ic{\{K => V\}} is a dynamic map from hashable keys \ty{K} to one
value type \ty{V}. Reading \ic{m[k]} always yields \ic{V?}; see
Section~\ref{sec:maps}.

\ty{K} is a hashable key type (\ic{str}, \ic{long}, \ic{byte},
\ic{bool}, \ic{float}), but a \textbf{map literal only accepts \ic{str}
keys} --- the checker's inference for a literal fixes the key type to
\ic{str}, so \ic{\{ 1: 5 \}} infers a shape rather than a map. Build a
non-\ic{str}-keyed map from an annotated empty literal and store into
it (Section~\ref{sec:maps}).

\subsection{Options and pointers}
\label{sec:type-options}

\ic{T?} is ``a \ty{T} or nothing''. Its empty case is \ic{none};
\ic{some}/\ic{unwrap}/\ic{is\_some} and the postfix unwrap \ic{x!}
(Section~\ref{sec:builtin-option}) manipulate it. There is no
\ic{null}: \ic{none} is the only empty value, and it exists only at
option types --- a plain \ic{str} or \ic{long} can never be empty, so
\ic{s == none} on a \ic{str} is a type error. On its own \ic{none} has
the type ``option of anything'' (printed \ic{none}) and adopts the
option type its context expects. A plain \ty{T} does not widen to
\ic{T?}; write \ic{some(v)}. \ic{*T} is a raw machine pointer; \ic{\&e}
yields one (Section~\ref{sec:pointers}).

\subsection{Void}

\ic{void} is the type of an expression that computes no value: a loop,
an \ic{if} without \ic{else}, a call to a void function
(Section~\ref{sec:void}), a bare \ic{return}. Void is not a value: it
cannot be bound, passed, printed or returned as one.

\paragraph{Disagreeing \ic{if} arms.} When the two arms of an
\ic{if}/\ic{else} have different types, the whole \ic{if} is
\ic{void} --- a statement. That is legal wherever a statement is, and
an error wherever a value is needed:

\begin{lstlisting}
if (ok) { 1 } else { "one" };            // fine: a statement
let v = if (ok) { 1 } else { 2 };        // long
let w = if (ok) { 1 } else { "one" };    // error: cannot bind a void value to 'w'
\end{lstlisting}

\subsection{Functions}

\ic{(T, T, ...) -> R}. Return types are inferred when omitted and
checked when written. A parameter annotation is optional: without one
the parameter takes the unknown type, which poisons quietly
(Section~\ref{sec:poison}) and is treated as a machine word by the
backend --- so annotate every parameter that carries a string, array,
record or option. Function values are first-class (they may be bound,
passed, stored and called); see Section~\ref{sec:functions}.

\subsection{Annotations and inference}

\begin{itemize}[nosep]
\item \ic{def} parameter annotations are \textbf{recommended}; a missing
  one takes the unknown type and never cascades a second diagnostic
  (Section~\ref{sec:poison}).
\item Return types are inferred from the body when omitted; when written,
  every \ic{return} and the body's trailing value must conform.
\item \ic{let} annotations are optional and checked; an annotated
  \emph{array literal} \emph{adopts} the annotated element type:
  \ic{let b: [byte] = [104, 105];} packs those elements one byte each.
  Adoption is restricted to literals --- assigning an existing
  \ic{[long]} to a \ic{[byte]} binding is a type error (the layouts
  differ).
\item An empty array literal \ic{[]} is legal unannotated: there is
  nothing to infer an element type from, so the element type is
  \emph{unknown} and it poisons quietly (Section~\ref{sec:poison}).
  Annotate it to pin the element type --- \ic{let b: [byte] = [];}.
  An unannotated \ic{\{\}} is the empty record; under a map annotation
  it is the empty map: \ic{let m: \{str => long\} = \{\};}.
\end{itemize}

\subsection{Conversions}
\label{sec:conversions}

The implicit conversions are:
\ic{long $\bs to$ float} (the only numeric widening),
\ic{byte $\bs to$ long}, \ic{byte $\bs to$ float}, and the narrowing
\ic{long $\bs to$ byte} at a \ic{byte} position (the value is masked to
eight bits: \ic{let b: byte = 258;} stores 2). A \ic{str} also
flows into a \ic{*T} position, since a string value is a byte address.
Apart from these, nothing converts: there is no cast syntax.

A value converts where it is stored under a known type: an annotated
\ic{let}; an assignment to a variable, element, field, map entry or
pointee; an argument; a \ic{return} or trailing value under a declared
return type; an element of an array literal; and the element that
\ic{+=}, \ic{*=} or \ic{::} adds. Array \ic{+} and \ic{+=} with an
array operand convert element by element into the left operand's
element type: \ic{[0.5] + [7, 8]} is \ic{[0.5, 7.0, 8.0]}. Arithmetic
on a \ic{byte} is \ic{long} arithmetic (\ic{-b} is a negative
\ic{long}); storing the result back into a \ic{byte} masks it.

Nothing converts a value that is already inside another one. Arrays,
maps and records are references and an option is a box, so an existing
\ic{[long]} is not a \ic{[float]}, and a record \ic{\{x: long\}} is not
a \ic{\{x: float\}}: a record's numeric field types must match exactly.

\subsection{Operator typing}

Arithmetic requires numbers; if either side is \ic{float} the result is
\ic{float} and the other side widens. \ic{+} additionally allows
strings (concatenation, with auto-stringification of the non-string
side) and arrays (Section~\ref{sec:arrays}); arrays win over strings
when both appear. Comparisons yield \ic{bool}; \ic{<}\,\ldots\ic{>=}
on two strings is lexicographic byte order. \ic{\&\&}/\ic{||} accept any
values: a non-zero value is true and the result is \ic{bool}
(Section~\ref{sec:evaluation}). Equality works on any two values of the
same type; it is by content for scalars and strings, structural for
arrays, and by identity for records and maps (Section~\ref{sec:equality}).

\subsection{Compound assignment typing}
\label{sec:opassign-typing}

For a target of type \ty{T} and value of type \ty{U}:

\begin{center}
\begin{tabular}{lll}
\toprule
\ty{T} & operators & rule \\
\midrule
numbers & \ic{+= -= *= /= \%=} & \ty{U} converts to \ty{T}; stored as \ty{T} \\
\ic{str} & \ic{+=} only & \ty{U} $=$ \ty{str} \\
array & \ic{+=}, \ic{*=} only & Section~\ref{sec:arrays} \\
\bottomrule
\end{tabular}
\end{center}

\noindent Anything else is a compile error, e.g. ``\ic{arrays support
only '+=' and '*=' (in place), not '-='}''. \ic{++}/\ic{--} require a
numeric target.

\subsection{The poison law}
\label{sec:poison}

Missing information (a missing annotation, an unanalyzable expression)
produces an \emph{unknown} type that unifies with everything. The law
this enables: \textbf{the root cause is reported exactly once} --- an
unknown-typed piece never generates follow-on errors at every enclosing
use. Concretely, unknown unifies silently in element checks, call
arguments and branches (so an unknown arm never makes an \ic{if}
void); the missing annotation itself is the single diagnostic.

\section{Expressions and evaluation}
\label{sec:evaluation}

Evaluation is strict and left-to-right; every operand is evaluated
exactly once per execution of the enclosing expression. Side effects
include calls, reads of places with side-effecting subexpressions, and
the mutating array operators.

\subsection{Numeric semantics}

\ic{long} arithmetic is two's complement modulo $2^{64}$: overflow
\emph{wraps}, it does not trap. Division and remainder truncate toward
zero (C semantics) and are subject to:

\begin{itemize}[nosep]
\item division or remainder by zero is a runtime error
  (Section~\ref{sec:errors});
\item \ic{LONG\_MIN / -1} evaluates to \ic{LONG\_MIN} (the quotient does
  not fit and would fault on x86) and \ic{LONG\_MIN \% -1} evaluates to
  \ic{0}. Implementations must special-case these two.
\end{itemize}

\ic{float} arithmetic is IEEE-754 binary64. Mixed long/float operands
widen the long side. Unary minus on a \ic{float} flips the sign, so
\ic{-0.0} is a negative zero (\ic{1.0 / -0.0} is $-\infty$) and
\ic{-(-x)} is \ic{x} for every \ic{x}.

\subsection{Logic}

\ic{\&\&} and \ic{||} short-circuit: the right operand is evaluated
only when the left does not decide the result. Operands are not
required to be \ic{bool} --- any non-zero value is true, and the result
of either operator is a \ic{bool}.

\subsection{Strings}

Strings are immutable values: concatenation \ic{a + b} produces a new
string. If exactly one side of \ic{+} is a string, the other side is
stringified (Section~\ref{sec:builtin-tostring}) and concatenated.
A format string, \ic{f"\{x\} of \{n:>4\}"}, writes such a concatenation
with fields (Sections~\ref{sec:literals} and~\ref{sec:format-spec}).
Comparisons \ic{<} \ic{<=} \ic{>} \ic{>=} on two strings compare bytes
lexicographically; \ic{==} compares contents.

\paragraph{Appending.}\label{sec:str-append} \ic{s += t} means \ic{s = s + t}, but a loop of
appends is linear, not quadratic: when \ic{s} is a variable whose
address is never taken and that no closure captures, and \ic{s} still
holds the string its own previous \ic{+=} built, the bytes of \ic{t} are
written in place behind it (the buffer doubles when full). Any other
use of \ic{s} --- a read, a copy into another binding, an element or
argument --- ends that ownership, and the next append copies once, so
an append never changes a string that anything else can see. Building
a string with \ic{+=} costs amortized $O(1)$ per appended byte, like
pushing onto an array.

\subsection{Assignment}
\label{sec:assignment}

\ic{p = e} evaluates \ic{e} and stores it into the place \ic{p}.
Assignment is a statement form, not an operator: it cannot appear inside
another expression, so \ic{let y = (x = 7);} and \ic{let y = x = 7;} are
syntax errors (so are the compound forms of
Section~\ref{sec:exactly-once}; only \ic{++} and \ic{--} nest). As the
last statement of a block it gives the block its value, the value
stored: \ic{let y = \{ x = 5 \};} sets both \ic{x} and \ic{y} to 5, and
a function body ending in an assignment returns the assigned value.
Section~\ref{sec:references}
defines what storing into arrays, records and variables means for
aliasing.

\subsection{Increment, decrement and compound assignment}
\label{sec:exactly-once}

The four increment/decrement forms and the five compound-assignment
operators share one rule:

\begin{quote}
\textbf{Exactly-once rule.} The target place's subexpressions are
evaluated \emph{exactly once}. Writing \ic{a[f()] += e} evaluates
\ic{f()} one time, no matter how the place is read and written.
\end{quote}

\noindent Formally, let a \ty{place} be $p ::= x \mid a[i] \mid o.f$,
possibly nested. Define $\mathrm{cap}(p)$, the capture of a place into
fresh variables (every non-name subexpression runs once):

\begin{align*}
\mathrm{cap}(x) &= \epsilon \\
\mathrm{cap}(a[i]) &= \mathrm{cap}(a);
   \;\; \text{if } a \text{ is not a name: } \texttt{let } \beta = a \texttt{;}
   \;\; \texttt{let } \iota = i \texttt{;} \\
\mathrm{cap}(o.f) &= \mathrm{cap}(o);
   \;\; \text{if } o \text{ is not a name: } \texttt{let } \beta = o \texttt{;}
\end{align*}

\noindent with $\mathrm{read}(p)$ and $\mathrm{write}(p, v)$ the
corresponding read and store through the captured names. Then the
\emph{numeric} forms desugar as:

\begin{align*}
\texttt{p++} &\;\equiv\; \mathrm{cap}(p);\;\texttt{\{}~\texttt{let t = read}(p)\texttt{;}
   ~~\texttt{write}(p, \texttt{t + 1})\texttt{;}~~\texttt{t}~\texttt{\}} \\
\texttt{++p} &\;\equiv\; \mathrm{cap}(p);\;\texttt{\{}~\texttt{let t = read}(p)\texttt{ + 1;}
   ~~\texttt{write}(p, \texttt{t})\texttt{;}~~\texttt{t}~\texttt{\}} \\
\texttt{p--} &\;\equiv\; \mathrm{cap}(p);\;\texttt{\{}~\texttt{let t = read}(p)\texttt{;}
   ~~\texttt{write}(p, \texttt{t - 1})\texttt{;}~~\texttt{t}~\texttt{\}} \\
\texttt{--p} &\;\equiv\; \mathrm{cap}(p);\;\texttt{\{}~\texttt{let t = read}(p)\texttt{ - 1;}
   ~~\texttt{write}(p, \texttt{t})\texttt{;}~~\texttt{t}~\texttt{\}} \\[4pt]
\texttt{p op= e} &\;\equiv\; \mathrm{cap}(p);\;\texttt{\{}~\texttt{let t = read}(p)
   \texttt{ } op \texttt{ } e\texttt{;}~~\texttt{write}(p, \texttt{t})\texttt{;}~~\texttt{t}~\texttt{\}}
\end{align*}

\noindent for $op \in \{+,-,\times,\div,\%\}$ on numbers. The four
increment and decrement forms are expressions: postfix yields the
\emph{old} value, prefix the \emph{new} one. A compound assignment is,
like \ic{=}, a statement form (Section~\ref{sec:assignment}): it cannot
nest inside another expression, and as the last statement of a block it
yields the new value of the place (\ic{\{ x += 2 \}}).

\paragraph{Examples.}
\begin{lstlisting}
a[f()]++;         // { let y_t = f(); let v = a[y_t];
                  //   a[y_t] = v + 1; v }   -- f() runs once
rec.v[g()] += 1;  // g() once; the read and the write use the same element
x.y.z += e;       // x and x.y evaluated once
\end{lstlisting}

\paragraph{Arrays.} On array targets \ic{+=} and \ic{*=} are
\emph{primitive in-place} operations --- they do not desugar through
\ic{+} (which copies) --- and are defined in Section~\ref{sec:arrays}.
On string targets only \ic{+=} is legal and concatenates (in place when
it can, Section~\ref{sec:str-append}). Every other
combination is a type error (Section~\ref{sec:opassign-typing}).

\paragraph{Maps and slices.} A map store \ic{m[k] = v} evaluates \ic{m},
\ic{k} and \ic{v} once each; a slice assignment \ic{a[lo:hi] = b} the same
way (Section~\ref{sec:arrays}).

\section{Arrays}
\label{sec:arrays}

Arrays are the language's one growable sequence. They are
\textbf{reference values}: a binding names an array \emph{object}, and
mutating operators write into that object so every binding sees the
change. Arrays are \textbf{homogeneous}: one element type, enforced
everywhere.

\subsection{Structure: length and capacity}
\label{sec:lencap}

An array object carries two independent quantities:

\begin{center}
\begin{tabular}{lll}
\toprule
 & meaning & invariant \\
\midrule
\ic{len} & number of live elements & $\texttt{len} \ge 0$ \\
\ic{cap} & elements the storage can hold & $\texttt{len} \le \texttt{cap}$ \\
\bottomrule
\end{tabular}
\end{center}

\noindent Indexing is trapped: reading or writing \ic{a[i]} with
\ic{i < 0} or \ic{i >= len} is a runtime error (``index out of
bounds''). \ic{cap} is never observable through operators --- only
\ic{len} is (\ic{len(a)}).

\paragraph{Growth policy (amortized doubling).}
\begin{align*}
\texttt{cap}_0 &= 0 \\
\texttt{cap}_{k+1} &= \max(8,\; 2\cdot\texttt{cap}_k)
   \qquad \text{when } \texttt{len} \text{ reaches } \texttt{cap}
\end{align*}

\noindent That is: the array grows \emph{only} when the capacity is
reached; the first growth yields capacity 8 and every later growth
doubles it ($8, 16, 32, \ldots$). Appending $n$ elements one at a time
therefore costs $O(n)$ copies total, $O(1)$ amortized per element. The
capacity never shrinks.

\subsection{Operators}
\label{sec:array-ops}

Let \ty{a}, \ty{b} be arrays of \ty{T}, and \ty{x} a value of type
\ty{T}. The complete operator semantics:

% The four columns are wide; a tighter \tabcolsep keeps the table inside
% the text block.
\begingroup
\setlength{\tabcolsep}{3.4pt}
\begin{center}
\begin{tabular}{llll}
\toprule
expression & result & effect on operands & value types \\
\midrule
\ic{a + b} & new array $= a \mathbin{+\!\!+} b$ & none &
   \ty{T} $=$ \ty{T} \\
\ic{a + x} & new array $= a \mathbin{+\!\!+} [\,x\,]$ & none &
   type of \ty{x} is \ty{T} \\
\ic{x + a} & new array $= [\,x\,] \mathbin{+\!\!+} a$ & none &
   type of \ty{x} is \ty{T} \\
\ic{rows + row} & row appended whole & none & \ic{[[T]]}, \ic{[T]} \\
\ic{row + rows} & row prepended whole & none & \ic{[[T]]}, \ic{[T]} \\
\ic{x :: a} & new array $= [\,x\,] \mathbin{+\!\!+} a$ & none &
   \ty{x} is ONE element \\
\ic{a += b} & \ty{a} & appends every element of \ty{b}, in place &
   \ty{T} $=$ \ty{T} \\
\ic{a += x} & \ty{a} & pushes \ty{x}, in place &
   type of \ty{x} is \ty{T} \\
\ic{a *= x} & \ty{a} & prepends \ty{x}, in place &
   type of \ty{x} is \ty{T} \\
\ic{a = b}  & \ty{b} & aliases: both names name one array & \\
\bottomrule
\end{tabular}
\end{center}
\endgroup

\begin{sloppypar}
\noindent The two \ic{rows} rows are the array-as-an-element case, and
they say what \ic{[[T]]} plus a \ic{[T]} does: the array goes in
\emph{whole}, as one element, because a \ic{[T]} is what that array
holds. \ic{rows + rows} (both \ic{[[T]]}) concatenates as usual, and so
does \ic{rows += row}, which appends the row in place.
\end{sloppypar}

\noindent In prose, matching the design intent exactly:

\begin{itemize}[nosep]
\item \ic{arr1 + arr2} --- make a new array which is \ic{arr1}
  concatenated with \ic{arr2}; the operands are untouched.
\item \ic{arr1 += arr2} --- append \ic{arr2} to the end of \ic{arr1},
  \emph{in place}.
\item \ic{arr1 = arr2} --- both point to the same array in memory.
\item \ic{arr1 + elem} --- a new array which is \ic{arr1} concatenated
  with \ic{[elem]}; only allowed if \ic{elem} has, or converts to, the
  element type (Section~\ref{sec:conversions}).
\item \ic{elem + arr1} --- a new array which is \ic{[elem]} concatenated
  with \ic{arr1}.
\item \ic{+=} means in-place push (of an element) or append (of an
  array); \ic{+} always means a \emph{new} array; \ic{*=} means prepend
  the item in place. A push and a prepend both take amortized constant
  time: an array keeps spare room at both ends and doubles it when it
  runs out, so building an array from either end is linear.
\end{itemize}

\paragraph{The one-thing rule.} Whether \ic{+=} pushes or appends is
decided by the \emph{shape of the value against the target's element
type}: an array whose own elements are the target's element type has
those elements appended; anything else --- a single value, or an array of
some other element type --- is pushed as one element. So a \ic{[long]}
onto a \ic{[[long]]} is pushed whole (it is one \ic{[long]} element), a
\ic{[[long]]} onto a \ic{[[long]]} appends its rows, and a \ic{[str]} is
neither, so it is an error:

\begin{lstlisting}
let rows = [[1], [2]];
rows += [3];        // pushes the one row [3] -- three rows now
rows += [[4], [5]]; // appends two rows -- five rows now
rows += 6;          // error: '+=' element: expected [long], got long
rows += ["s"];      // error: '+=' appended elements: expected [long], got str
\end{lstlisting}

\paragraph{Homogeneity.} Every mixing form requires matching element
types; the checker rejects mismatches --- including the append case, where
a non-matching array would splice wrong-typed elements:

\begin{lstlisting}
let xs = [1, 2];
let fs: [float] = [0.5];
xs += fs;           // error: '+=' appended elements: expected long, got float
\end{lstlisting}

\paragraph{\ic{a += a}.} The right side is captured (its length snapshotted)
before appending, so \ic{a += a} doubles \ic{a} with a copy of its
original elements.

\subsection{Aliasing and visibility}
\label{sec:array-aliasing}

Assigning arrays aliases; mutating operators are visible through every
alias:

\begin{lstlisting}
let m1 = [1, 2];
let m2 = m1;      // ONE array, two names
m1 += 3;          // m2 sees it: len(m2) == 3
\end{lstlisting}

\noindent By contrast \ic{+} never mutates: \ic{let ys = xs + [9];}
leaves \ic{xs} at its old length.

\subsection{Fixed arrays}

\ic{[T; N]} arrays are allocated with exactly \ic{N} slots; their length
is part of the type, so they reject growth: \ic{+=} and \ic{*=} on a
fixed array is a compile error (``cannot grow fixed array \ic{[long; 3]}
with \ic{'+='}''). Indexing is bounds-trapped like heap arrays. A fixed
array may be sliced (Section~\ref{sec:slices}) and slice-assigned
(below); the result of a slice is a heap \ic{[T]}.

Since the length is in the type, an annotated fixed array checks what it
is given: a heap \ic{[T]} never converts to a \ic{[T; N]}, and a literal
is checked against the size it claims --- \ic{let v: [long; 5] = [1, 2];}
is an error (\ic{expected [long; 5], got [long]}), while the matching
\ic{[long; 2]} is fine. Build one with \ic{[v; n]}, which is a fixed array
of \ty{n} copies of \ic{v} (Section~\ref{sec:expr-grammar}).

\subsection{Slices}
\label{sec:slices}

\ic{a[lo:hi]} evaluates to a \textbf{new} heap \ic{[T]} holding the
elements in $[\texttt{lo}, \texttt{hi})$. Missing bounds default to
\ic{0} and \ic{len}; bounds are clamped (no trap). All four spellings
are legal: \ic{a[:]}, \ic{a[lo:]}, \ic{a[:hi]}, \ic{a[lo:hi]}. Strings
slice the same way: \ic{s[lo:hi]} is a new \ic{str} of the bytes in that
range, clamped alike. There is no \ic{substr}.

\begin{sloppypar}
\paragraph{Slice assignment.} \ic{a[lo:hi] = b} writes \ic{b} into \ic{a}
in place, over the range \ic{[lo, hi)} --- the same clamped range a read
would take, on either kind of array. It is a statement, so as the last
one of a block it is that block's value (Section~\ref{sec:assignment}).
The target's subexpressions, the bounds and \ic{b} are each evaluated
once, left to right (Section~\ref{sec:exactly-once}). The right side
must be an array whose element type is the target's, a \ic{str} cannot
be slice-assigned, and the length is checked when the program runs: a
\ic{b} of another length traps with
\ic{runtime error: slice assignment length mismatch}, and a range
outside the array with \ic{runtime error: slice assignment out of
bounds}. The static failures name themselves:

\begingroup
\setlength{\tabcolsep}{4pt}
\begin{center}
\begin{tabular}{ll}
\toprule
mistake & message \\
\midrule
a \ic{[float]} into a \ic{[long]} &
  \ic{slice assignment: cannot store [float] into [long]} \\
not an array &
  \ic{slice assignment needs an array on the right, got long} \\
a \ic{str} target &
  \ic{cannot assign to a slice of str (strings are immutable)} \\
\bottomrule
\end{tabular}
\end{center}
\endgroup
\end{sloppypar}

\begin{lstlisting}
let a = [1, 2, 3, 4, 5];
a[1:3] = [9, 8];            // a is [1, 9, 8, 4, 5]
a[1:3] = [7];               // traps: slice assignment length mismatch
\end{lstlisting}

\subsection{Byte strings}

A \ic{[byte]} array \emph{is} a binary string: \ic{b"..."} literals
produce them, \ic{+} concatenates, slices cut, and printing one shows
the Python-style repr --- printable ASCII verbatim, quote and backslash
escaped, \ic{\bs n} \ic{\bs r} \ic{\bs t} as escapes, and every other
byte as \ic{\bs xHH}:

\begin{lstlisting}
print(b"\bs x7fELF");     // b'\bs x7fELF'
let hdr = b"AB" + b"\bs x00";
print(hdr);              // b'AB\bs x00'
print(hdr[0:2]);         // b'AB'
\end{lstlisting}

\noindent \ic{\bs{std/struct.y}} builds on this with pack/unpack
helpers for file headers and binary layouts (\ic{u32be}, \ic{get\_i16le},
\ic{pack}, \ic{unpack}, \ldots).

\subsection{Operators are the interface}

There are no array mutator functions: \ic{+=} and \ic{*=} \emph{are} push
and prepend, and there is no \ic{push}, \ic{prepend}, \ic{map\_get} or
\ic{map\_set} to call instead (Section~\ref{sec:builtin-pointers}).
\ic{len(a)} is the only builtin whose argument is an
array \emph{as such}; \ic{join(xs, sep)} also takes an array but of any
element type (each element is stringified), and \ic{bytes}/\ic{from\_bytes}
and \ic{fs.write\_bytes} take a \ic{[byte]}.

\section{Records and maps}
\label{sec:references}

Y has \emph{two} key/value types, deliberately distinct:

\begin{center}
\begin{tabular}{lll}
\toprule
 & \textbf{record} (static) & \textbf{map} (dynamic) \\
\midrule
type & \ic{\{ a: long, b: str \}} & \ic{\{str: long\}} \\
keys & identifiers, fixed at compile time & any hashable value at run time \\
values & may differ per field & one homogeneous type \ty{V} \\
layout & C struct: fields at fixed offsets & growable key/value store \\
keying & \ic{r.a} & \ic{m['a']} \\
\bottomrule
\end{tabular}
\end{center}

\subsection{The value model}

\begin{center}
\begin{tabular}{lll}
\toprule
category & types & assignment \\
\midrule
values & \ic{long} \ic{byte} \ic{bool} \ic{float} \ic{str} \ic{null}
   & copies the bits \\
references & arrays (\ic{[T]}, \ic{[T; N]}), records, maps
   & \textbf{aliases} the object \\
\bottomrule
\end{tabular}
\end{center}

\noindent ``Reference'' means: every binding of the same object names
one storage, and every in-place mutation (\ic{+=}, \ic{*=}, field
assignment, map stores) is observable through all of them.

\subsection{Records}

A record literal \ic{\{ name: "y", n: 1 \}} allocates one record object
whose shape is the literal's fields \emph{as written}; field order is
preserved. At run time a record is exactly a C struct: field \ty{i} is
the word at \ic{+8i} from the record address, so \ic{r.a} compiles to a
single load at a compile-time offset.

\begin{lstlisting}
let r = { v: [7], n: 1 };
let s = r;            // aliases the record
r.v += 8;             // s.v sees it
s.n = 2;              // r.n sees it
\end{lstlisting}

\noindent The shape is fixed: writing a field the shape does not have is
a \emph{compile error} (there is no runtime field growth), and a field
read requires the field to exist statically.

\paragraph{Record spread.} \ic{\{a: 0, ...obj\}} builds a record whose
shape is the \emph{union} of the two shapes — entirely statically; no
dynamic shape is involved. Entries merge with last-wins values and
first-appearance positions, so \ic{obj} may overwrite \ic{a} and a later
explicit key may overwrite what \ic{obj} spread:

\begin{lstlisting}
let b = {x: 1, y: 2};
let c = {y: 9, z: 3};
print({a: 0, ...b});     // {a: 0, x: 1, y: 2}
print({...b, ...c});     // {x: 1, y: 9, z: 3}
print({...b, y: 100});   // {x: 1, y: 100}
\end{lstlisting}

\subsection{Maps}
\label{sec:maps}

A map type is \ic{\{\ty{K}: \ty{V}\}} where \ty{K} is a hashable key
type (\ic{str}, \ic{long}, \ic{byte}, \ic{bool}, \ic{float}) and \ty{V}
is one homogeneous value type. Map literals use quoted keys (identifier
keys are record syntax):

\begin{lstlisting}
let m = {"a": 1, "b": 2};      // {str: long}
let empty: {str: long} = {};   // empty literals need the annotation
m["c"] = 3;                    // insert or overwrite
print(m["a"]);                 // 1
print(m["zz"]);                // null -- reads yield V?
\end{lstlisting}

\noindent Keying is symbolic and exclusive: \ic{m['k']} is the map form
and is reserved for maps; \ic{r.a} is the record form. Reading
\ic{m[k]} yields \ic{\ty{V}?} — \ic{null} when the key is absent — and
\ic{m[k] = v} stores in place (the store is visible through every
binding of the map). \ic{len(m)} is the number of entries. The map grows
as needed (amortized doubling, like arrays).

\subsection{Equality}

\ic{==} on two arrays or two records is \textbf{structural}: same
length/fields, pairwise-equal contents. It is never a pointer
comparison; \ic{arr == []} is an emptiness test.

\subsection{Copying}

There is no implicit deep copy. To duplicate an array use \ic{+} with an
empty right side (\ic{a + []}) or slicing (\ic{a[:]}); records can be
rebuilt field by field. The language deliberately makes sharing the
default and copying explicit.

\section{Functions, closures, options, pointers}
\label{sec:functions}

\subsection{Functions}

A \ic{def} declares a function; parameters are typed, the return type is
inferred or checked. Write recursion with \ic{self} --- that is the
portable spelling, and it is what a function \emph{literal} (which has no
name) has to use. A \emph{named} \ic{def}'s own name is also in scope
inside its body, so \ic{count(n - 1, acc + 1)} and \ic{self(n - 1, ...)}
are the same call:

\begin{lstlisting}
def fact(n: long) -> long {
    if (n <= 1) { 1 } else { n * self(n - 1) }
};
\end{lstlisting}

\noindent Functions are first-class values: they may be bound
(\ic{let f = g;}), stored in arrays and records, passed to other
functions, and called through any of those (\ic{f(1)}, \ic{fs[0](1)},
\ic{r.handler(1)}). Calling a non-function is a compile error.

Function types include both parameter and return types. A function may also
be passed where the same parameters and a \ic{void} result are expected: the
caller discards its result. Other result types must agree.

\subsection{Tail calls}

A self-call in \emph{tail position} --- the body's trailing expression,
a \ic{return}'s value, or an if-arm's trailing expression --- recycles
the stack frame instead of growing it: the arguments are evaluated, the
frame is torn down, and control re-enters the function. Tail recursion
therefore runs in constant stack space:

\begin{lstlisting}
def sum_to(n: long, acc: long) -> long {
    if (n == 0) { acc } else { self(n - 1, acc + n) }
};
\end{lstlisting}

\noindent Calls outside tail position (\ic{self(x) + self(y)},
\ic{let t = self(x)}, inside loops) are ordinary calls and stack as
usual.

\subsection{Closures}
\label{sec:closures}

A function literal captures the enclosing bindings it reads --- and only
those that are already in scope where the literal is written (a name
bound by a \ic{let} \emph{after} the \ic{def} is not in scope and is an
undefined-variable error, exactly like reading a local before its
declaration). Captured bindings are \textbf{shared}: a write inside the
closure is visible outside, and a write outside is visible inside.

A literal written as an \emph{arrow} is the same closure with less
punctuation: \ic{v => v + 1} is \ic{def(v) \{ v + 1 \}},
\ic{(e, i) => e + i} takes two parameters, and \ic{() => 0} takes none.
The body is one expression, or a block; a parameter may carry a type
(\ic{(acc: long, v: long) => acc + v}), and the parentheses are optional
for a lone parameter. Arrow literals appear anywhere a \ic{def} literal
does, and they are what the array methods take
(Section~\ref{sec:array-methods}).

\begin{lstlisting}
def counter() {
    let n = 0;
    def tick() -> long { n = n + 1; n };
    tick
};
let t = counter();
t(); t();
print(t());        // 3 -- the one shared n
\end{lstlisting}

\noindent Each \emph{invocation} of the enclosing function creates fresh
bindings, so closures created in different iterations see different
cells (``\ic{let} semantics''):

\begin{lstlisting}
while (i < 3) {
    let x = i;
    fns += def() -> long { x };
    i = i + 1;
};
\end{lstlisting}

\noindent The \emph{loop variable} of a \ic{for x in a..b} or a
\ic{foreach (x in a)} is fresh on each iteration, like the \ic{let}
inside the loop above --- it is the one binding the loop makes per
iteration rather than one the body does:

\begin{lstlisting}
let fns: [() -> long] = [];
for i in 0..3 {
    fns += def() -> long { i };     // this iteration's i: 0, 1, 2
};
foreach (e in [7, 8, 9]) {
    fns += def() -> long { e };     // 7, 8, 9
};
\end{lstlisting}

\noindent A binding that is \emph{not} a loop variable --- one declared
before the loop, or the \ic{let} inside its body --- is a single cell,
and every closure the loop makes shares it, so they all see where the
loop ended:

\begin{lstlisting}
let k = 0;
while (k < 3) {
    k = k + 1;
    fns += def() -> long { k };     // one k, shared: every call says 3
};
\end{lstlisting}

\subsection{Options}
\label{sec:builtin-option}

\ic{T?} is the type of ``a \ty{T} or nothing''; its empty case is
\ic{none}. \ic{some(v)} boxes a value as present, \ic{is\_some(v)}
tests, and \ic{unwrap(v)} --- or its postfix spelling \ic{v!} ---
returns the payload and is a \emph{runtime error} on \ic{none} (it is a
programmer error, loudly). \ic{x == none} is the is-empty check, and
two options are equal when both are \ic{none} or both hold equal payloads (Section~\ref{sec:equality});
an option of an array is an optional reference, not a copied one.

\begin{lstlisting}
let line = readline();          // str?
if (line == none) { print("eof"); } else { print(line! + "!"); };
\end{lstlisting}

\subsection{Void}
\label{sec:void}

A \ic{def} whose result is not a value is \ic{void}. Write \ic{-> void}
to say so: the body's trailing value is then discarded, and \ic{return}
takes no value. An unannotated \ic{def} whose body ends in a statement
infers \ic{void}, and then every \ic{return} in it must be bare: a
\ic{return 5;} there is an error, and so is a bare \ic{return;} in a body
that ends in a value. A void call is a statement; using it as a value --- binding
it, printing it, passing it as an argument --- is a compile error:

\begin{lstlisting}
def log(msg: str) -> void { print("log: " + msg); };
log("hi");                 // fine: a statement
let x = log("hi");         // error: cannot bind a void value to 'x'
\end{lstlisting}

\subsection{Pointers}
\label{sec:pointers}

\ic{*T} is a raw address of one \ty{T}; \ic{*p = v} stores,
\ic{*p} reads. The address-of operator \ic{\&e} is deliberately narrow:

\begin{itemize}[nosep]
\item \ic{\&x} --- address of a local or parameter;
\item \ic{\&a[i]} --- address of an array element (any array expression:
  \ic{\&grid[0][i]} is legal);
\item \ic{\&arr} on a whole array --- \emph{both} kinds, \ic{[T]} and
  \ic{[T; N]} --- is a compile error
  (``cannot take the address of a whole array``), and \ic{\&} of any
  temporary is a different one (``'\&' needs a plain variable``). A binding names
  the array \emph{header}, not its first element, so \ic{\&arr} would hand
  out the header; write \ic{\&arr[0]}, which works for either kind.
\end{itemize}

\noindent Pointers are for interfacing with the system-call level
(buffers for \ic{read}/\ic{write}-style calls); ordinary code uses
arrays and records. They support C-style arithmetic: \ic{p + n} and
\ic{p - n} move by whole \ty{T} elements (one byte for \ic{*byte}, the
width of a sized integer for \ic{*u32} and its kind, eight bytes for every
other pointer), \ic{n + p} is the mirror,
\ic{p - q} counts elements when both pointers have the same type, and
\ic{+=}, \ic{-=}, \ic{++} and \ic{--} scale the same way. Ordering
comparisons compare raw addresses; equality is word equality. The
offset must be \ic{long} or \ic{byte}, and mixing pointer types in a
difference is a compile error.

\noindent\textbf{Sized integers.} To read or write memory laid out by
someone else (a system-call structure, a file or wire format), a pointer's
target can be a sized integer: \ic{i8}, \ic{i16}, \ic{u16}, \ic{i32},
\ic{u32} and \ic{u64}, and the 16-, 32- and 64-bit ones with a \ic{be}
suffix for big-endian (\ic{*u32be}); \ic{le} spells the native order and
is the same type. (\ic{u8} is \ty{byte} and \ic{i64} is \ty{long}.) A sized
integer is \emph{only} a pointer's target: \ic{*p} is a \ty{long}, sign- or
zero-extended from the width in memory, and \ic{*p = v} stores the low
bytes of \ic{v} and drops the rest, as a store to a \ic{byte} does. A
variable, parameter, field or array element of type \ic{u32} is a compile
error (``u32 is only a pointer's target (write *u32)''): the width is a
property of the memory, not of the value. The pointer comes from
\ic{ptr\_from}, below, or from any \ic{*T} of the same type.

\noindent\textbf{Packed records.} A record normally keeps each field in a
word of its own, sorted by name; \ic{packed} lays the fields out in the order
written, byte after byte with no padding, so a record can describe a region of
memory someone else defined:
\begin{verbatim}
type Flags = bits u8 { ready, error, mode: 3, _: 2 };
type Hdr = packed { magic: u32be, len: u16, flags: Flags, tag: [byte; 4] };
let h: Hdr = { magic: 0xCAFEBABE, len: 513, flags: { ready: true, mode: 5 } };
h.len = h.len + 1;
h.flags.error = true;
let p: *Hdr = ptr_from(addr);   // fields over memory you did not make
\end{verbatim}
\noindent Inside \ic{packed} the sized integers of the last paragraph are field
types (\ic{u32be}, \ic{i16}, \ic{u8}\ldots), as are \ic{long}, \ic{float},
\ic{bool} (one byte), a pointer, a fixed array \ic{[T; N]} of those, and another
packed record, which sits inline. Anything with a header of its own (\ic{str},
\ic{[T]}, a map, an ordinary record) is refused when the type is used: the
collector does not read raw memory, so a reference stored in it would not be
kept alive. \ic{bits u8 \{ a, b, mode: 3 \}} (backing \ic{u8}, \ic{u16},
\ic{u32} or \ic{u64}) gives each name the next bits from the least significant
end: one bit, read as a \ty{bool}, unless a width follows; \ic{\_} skips bits.
Writing one bit field reads and rewrites the whole backing integer.

\noindent A packed value \emph{is} the address of its bytes. \ic{h.len} reads
or writes the field in place (widened by its sign, byte-swapped for \ic{be});
\ic{h.flags} and \ic{h.sub} are views of the nested bytes, so they take no
assignment as a whole, only their fields; an array field reads as a \ic{*T} to its
first element. A literal \ic{\{ .. \}} where a packed type is expected allocates
a zeroed block of \ic{sizeof(Hdr)} bytes and stores the fields it names.
\ic{Hdr} and \ic{*Hdr} convert to each other, and \ic{ptr\_addr(h)} is the
address. \ic{sizeof(T)} and \ic{offsetof(T, field)} are the layout as constants.
A packed value is compared and copied as an address, and cannot be printed.

\noindent\textbf{Pointers and integers.} A pointer is its address at run
time, and two builtins say so. \ic{ptr\_addr(p)} is the \ic{long} value of a
pointer. \ic{ptr\_from(n)} is the pointer whose address is the integer
\ic{n} (a \ic{long} or a \ic{byte}), the inverse. Its result has no pointee
type of its own, so the type comes from where it is used: write
\ic{let p: *byte = ptr\_from(a)}, pass it to a \ic{*T} parameter, or store it
in a \ic{*T} field. A bare \ic{let p = ptr\_from(a)} is a compile error
(``ptr\_from does not say which pointer type it makes''), and so is a
non-integer argument. Nothing checks that \ic{n} is mapped memory: a wrong
address faults when it is dereferenced, which is what a kernel's physical
memory access needs and why ordinary code does not use it.

\noindent\textbf{A pointer does not outlive what it points into.}
\ic{\&x} on a local or a parameter is the address of a cell in that
function's frame, and the frame is reclaimed the moment the function
returns: a pointer to a local that escapes reads whatever the stack
holds by then --- a number that was never written, or an address, and
it changes from run to run. Returning one is not a type error,
because the checker cannot see how long a value lives; the rule is that
a pointer into a frame may only be used while that frame is live, which
means in practice that it must not be returned and must not be stored
in anything outliving the call. Taking the address of a \emph{global}
is fine (its storage is permanent), and so is \ic{\&a[i]}: a heap
array's elements and a fixed array's live in a heap cell, never in a
frame. To hand data to a function or keep it, return the value, an
array or a record --- those never dangle --- and take a pointer to
\ic{\&a[i]} where the callee needs an address.

\section{Builtins}
\label{sec:builtins}

The following functions are available without declaration. ``Trap'' marks
a runtime error (Section~\ref{sec:errors}).

\subsection{Sequences and text}

\begin{description}[style=nextline,leftmargin=5.2em]
\item[\ic{len(x)}] Length of a string (bytes), array (elements) or map
  (entries) as \ic{long}. Records are rejected (a record's field count
  is part of its type).
\item[\ic{range(n)}] \ic{[0, 1, ..., n-1]}.
\item[\ic{trim(s)}] Strip leading/trailing space, tab, CR, LF.
\item[\ic{join(xs, sep)}] Concatenate \ic{[str]} with a separator.
\item[\ic{split(s, sep)}] Split on every occurrence of a non-empty
  separator; returns \ic{[str]} (empty for empty separator).
\item[\ic{index\_of(s, sub)}] First index of \ic{sub} in \ic{s}, or
  \ic{-1}.
\item[\ic{last\_index\_of(s, sub)}] Last index, or \ic{-1}.
\item[\ic{ord(s)}] First byte as \ic{long}, or \ic{-1} for empty.
\item[\ic{chr(n)}] One-byte string of \ic{n} (masked to 8 bits).
\item[\ic{byte\_at(s, i)}] Byte at \ic{i} as \ic{long}; traps out of
  bounds. Cheap per-byte access (no copy).
\end{description}

\noindent There is no \ic{substr}: a substring is a slice,
\ic{s[i:i + n]} (Section~\ref{sec:slices}).

\subsection{String methods}
\label{sec:string-methods}

A \ic{str} value takes methods: \ic{"str 1 2 3".split(" ")} is
\ic{["str", "1", "2", "3"]}. A method call is checked like a function
call (argument count and types). Method names are not reserved: a
function, a variable or a record field may share one, and \ic{r.f(x)} on
a record still calls the field. \ic{s.length} is a property, so
\ic{s.length()} is an error. Widths and positions count bytes.

\begin{description}[style=nextline,leftmargin=5.2em]
\item[Builtins as methods] \ic{s.split(sep)}, \ic{s.trim()} (also
  \ic{s.strip()}), \ic{s.index\_of(sub)}, \ic{s.last\_index\_of(sub)},
  \ic{s.ord()}, \ic{s.byte\_at(i)}, \ic{s.tonum()} and \ic{s.bytes()} are
  the builtins of the same name with \ic{s} as the first argument.
\item[\ic{sep.join(xs)}] \ic{join(xs, sep)}: the receiver is the
  separator.
\item[\ic{s.trim\_start()}, \ic{s.trim\_end()}] Strip space, tab, CR and
  LF at one end only (also spelled \ic{lstrip}, \ic{rstrip}).
\item[\ic{s.pad\_left(w, fill)}, \ic{s.pad\_right(w, fill)}, \ic{s.center(w, fill)}]
  Widen \ic{s} to \ic{w} bytes with \ic{fill} (default \ic{" "}, repeated
  and cut to fit); a longer \ic{s} is returned unchanged. \ic{center}
  puts the odd byte on the right.
\item[\ic{s.starts\_with(p)}, \ic{s.ends\_with(p)}, \ic{s.contains(p)}]
  \ic{bool}; \ic{""} is a prefix, suffix and part of every string.
\item[\ic{s.count(p)}] Non-overlapping occurrences, left to right
  (\ic{0} for an empty \ic{p}).
\item[\ic{s.replace(p, r)}] Every non-overlapping occurrence of \ic{p},
  left to right, replaced by \ic{r}; an empty \ic{p} replaces nothing.
\item[\ic{s.upper()}, \ic{s.lower()}] ASCII case mapping; other bytes are
  kept.
\item[\ic{s.repeat(n)}] \ic{s} written \ic{n} times; \ic{""} when
  \ic{n <= 0}.
\item[\ic{s.lines()}] Split at each LF, dropping a CR just before it; a
  final newline does not start an empty line, so \ic{"a\bs nb\bs n".lines()}
  is \ic{["a", "b"]} and \ic{"".lines()} is \ic{[]}.
\end{description}

\subsection{Conversions}
\label{sec:builtin-tostring}

\begin{description}[style=nextline,leftmargin=5.2em]
\item[\ic{tostring(v)}] String form: numbers as decimal (floats with six
  decimals), \ic{true}/\ic{false}, strings unchanged, options as their
  payload or \ic{"none"}, arrays as \ic{"[e0, e1, ...]"}, records as
  \ic{"\{k: v, ...\}"}, maps as \ic{"\{"k": v, ...\}"} (string keys
  quoted).
\item[\ic{tonum(s)}] Parse a decimal integer (\ic{0} on garbage).
\item[\ic{typeof(v)}] The static type as a printable value
  (\ic{typeof(none)} is \ic{"none"}).
\end{description}

\subsection{Binary data}
\label{sec:builtin-bytes}

\begin{description}[style=nextline,leftmargin=5.2em]
\item[\ic{bytes(s)}] \ic{[byte]} of the string's raw bytes.
\item[\ic{from\_bytes(b)}] String from a \ic{[byte]} or fixed \ic{[byte; N]} array.
\end{description}

\subsection{Options}
\label{sec:builtin-option-b}

\begin{description}[style=nextline,leftmargin=5.2em]
\item[\ic{some(v)}] A present option holding \ic{v}.
\item[\ic{is\_some(v)}] \ic{true} unless \ic{v} is \ic{none}.
\item[\ic{unwrap(v)}] The payload; \textbf{traps} on \ic{none}. The
  postfix form \ic{v!} is the same operation.
\end{description}

\subsection{Input/output and the environment}

\begin{description}[style=nextline,leftmargin=5.2em]
\item[\ic{print(x)}] Print a value followed by a newline to stdout.
\item[\ic{puts(s)}] Write a string to stdout, no newline added.
\item[\ic{eputs(s)}] Same, to stderr (keeps \ic{--time} output and
  program output apart).
\item[\ic{readline()}] One line from stdin without its newline, or
  \ic{none} at end of input (\ic{str?}). Stdin is buffered (16\,KB per
  \ic{read(2)}); before \ic{exec()} starts a child, the unread part is
  given back when stdin is a file, so the child reads on from the next
  line. On a pipe, lines already buffered stay with the program, as with
  C stdio.
\item[\ic{clear\_eof()}] No-op: \ic{readline()} reads again whenever
  its buffer is empty, so end of input is never latched. Present so the
  same source runs on every compiler.
\item[\ic{read\_file(path)}] File contents as \ic{str?}: \ic{none} when
  the file cannot be read.
\item[\ic{write\_file(path, s)}] Write a string; \ic{bool} success.
\item[\ic{write\_bytes(path, b, mode)}] Write \ic{[byte]} binary-safe
  and chmod (\ic{mode == -1} skips chmod --- any other value is passed
  to \ic{fchmodat} verbatim, so \ic{493} is \ic{0755}); \ic{bool} success.
\item[\ic{argv()}] \ic{[str]} of the program's arguments.
\item[\ic{exit(code)}] Terminate with an exit code.
\item[\ic{clock\_us()}] Monotonic microseconds as \ic{long}.
\item[\ic{heap\_used()}] Bytes allocated so far.
\item[\ic{sleep(ms)}] Suspend the process (\ic{void}).
\item[\ic{syscall(...)}] Raw Linux syscall wrapper; the first argument
  is the syscall number.
\item[\ic{exec(cmd)}] Run a shell command, return its exit code.
  Removed from the binary entirely by \ic{--no-exec} (then calling it
  is a compile error).
\end{description}

\subsection{Pointers}
\label{sec:builtin-map}

\begin{description}[style=nextline,leftmargin=5.2em]
\item[\ic{ptr\_addr(p)}] The address of a pointer as \ic{long}.
\end{description}

\noindent Mutating and keying are symbolic, not functional: arrays use
\ic{+=} / \ic{*=}, records use \ic{r.f}, maps use \ic{m['k']} and
\ic{m['k'] = v}. The legacy functions \ic{push}, \ic{prepend},
\ic{map\_get} and \ic{map\_set} do not exist.

\section{Errors and diagnostics}
\label{sec:errors}

\subsection{Compile time}

Static checking is fatal: a rejected program compiles to nothing.
Diagnostics are one line each,

\begin{center}
\ic{line:col: error: message}
\end{center}

\noindent positioned in the user's source even when the compilation
splices prelude/runtime code in front of it. The checker reports a root
cause \emph{once} (Section~\ref{sec:poison}): a missing annotation or
unanalyzable expression poisons quietly instead of cascading errors
through every enclosing use. Typical families:

\begin{itemize}[nosep]
\item name errors: \ic{undefined variable 'x'};
\item shape errors: \ic{record has no field 'f'};
\item operator typing: \ic{'+=' appended elements: expected [long], got
  long}, \ic{arrays support only '+=' and '*=' (in place), not '-='},
  \ic{'++' needs a number target};
\item call errors: \ic{call takes 2 args, got 1};
\item non-value use: \ic{cannot index long (only arrays and maps)};
\item void use: \ic{cannot bind a void value to 'x'}.
\end{itemize}

\subsection{Run time}

The compiled program traps --- writes a one-line message to stdout and
exits with status 1 --- on:

\begin{center}
\begin{tabular}{ll}
\toprule
condition & message \\
\midrule
array index outside \ic{[0, len)} & \ic{runtime error: index out of bounds} \\
division or remainder by zero & \ic{runtime error: division by zero} \\
\ic{unwrap(none)}, \ic{none!} & \ic{runtime error: unwrap of none} \\
missing dynamic field/key & \ic{runtime error: record has no field 'k'} \\
slice assignment out of range & \ic{runtime error: slice assignment out of bounds} \\
match with no arm covering a value & \ic{runtime error: no match arm matched} \\
memory exhausted & \ic{runtime error: out of memory} \\
unbounded (non-tail) recursion & stack exhaustion: no limit, so the
  process dies on the guard page \\
\bottomrule
\end{tabular}
\end{center}

\noindent These messages go to \textbf{stdout} (the trap stubs use
\ic{write(1, ...)} directly), so a program's real output stays
recognizable; \ic{eputs} is the way to send diagnostics to stderr.

\noindent Fixed arrays trap exactly like heap arrays; slices never trap
(they clamp). Integer overflow is \emph{not} an error
(Section~\ref{sec:evaluation}).

\subsection{Exit status}

A successfully run program exits with the code it chose
(\ic{exit(code)} or falling off the end, which is \ic{0}); a trapped
program exits with \ic{1}. The compiler itself exits \ic{0} on success,
\ic{1} on a rejected input.

\subsection{Reproducibility}

The compiler is deterministic: the same input produces byte-identical
output on every run and in either implementation. This is checked
continuously by comparing the two compilers' output binaries byte for
byte, and by self-compilation (a compiler built by the compiler is
identical to the one that built it).

\section{Y-OS and yunix}
\label{sec:os}

\textsf{Y} has a second compilation target beside the hosted Linux ELF:
\ic{ycc --target=kernel} compiles a program that boots on bare metal, and \emph{Y-OS}
is a small operating system built from such programs. It is a target, not a
mode on top of Linux: \ic{res\_target} becomes \ic{"kernel"} and everything
downstream that would otherwise assume an operating system changes behind one
predicate, \ic{target\_is\_kernel()}. A hosted compile is untouched
(Section~\ref{sec:os-cost}).

Y-OS is written in \textsf{Y} alone. The kernel proper is the standard units
\ic{kernel}, \ic{kmem}, \ic{yfs} and \ic{fd}, which travel in the std pack
like every other unit, so a lone \ic{ycc --target=kernel} in an empty directory
builds a bootable kernel with no checkout; the \emph{system} on top of it is
\ic{os/kernel/*.y}, the shell and its programs in \ic{os/uland/}, and the
host-side image builder \ic{os/tools/mkfs.y}.

\subsection{The kernel target}

\begin{lstlisting}
$ ./ycc --target=kernel prog.y -o prog.elf
$ qemu-system-x86_64 -kernel prog.elf -serial stdio -display none
\end{lstlisting}

\noindent The output is an ELF32 Multiboot-1 payload: the entry is 32-bit
code that builds paging, sets PAE/LME/PG and far-jumps into long mode,
because that is the mode a Multiboot loader starts an image in --- the one
step the in-tree assembler cannot express, whose encoder is x86-64 only. The
image is linked at 1\,MB, the serial port COM1 is its output, and \ic{print}
and \ic{puts} work there unchanged (\ic{y\_os\_write} is replaced by a COM1
writer). There is no libc, no Linux and no \ic{/bin}.

\paragraph{What a kernel can use.} The runtime is the same prelude with the
few blocks that reach an operating system replaced: the allocator bumps
through a kernel heap, \ic{y\_os\_read} blocks on the 16550, \ic{y\_os\_exit}
halts the machine, and the five trap stubs print and stop. The string, array
and float cores are pure computation and are used as they are, and the tree
shaker drops what the program does not call exactly as it does for a hosted
program.

A kernel refuses the std units that are one Linux system call away from
working --- \ic{fs}, \ic{net}, \ic{exec}, \ic{prelude}, \ic{crypto}, \ic{x509},
\ic{tls}, \ic{http}, \ic{tui}, \ic{args}, \ic{random}, \ic{test}, \ic{img},
\ic{audio}, \ic{ipc} and \ic{snd} --- each with a one-line reason
(\ic{kernel\_std\_why}), because ``not available'' without one reads as a
compiler limitation. The rule is a deny list, not an allow list: a unit that
is plain computation (\ic{json}, \ic{regex}, \ic{math}, \ic{struct}, \ic{utf8})
is available in a kernel the day it lands. The reverse rule is \ic{std/kernel},
which is \ic{ndef}s into assembly only a kernel image has, so a hosted compile
refuses \ic{import kernel} and says so. Both rules fire only when the unit
resolved \emph{is} the standard one: a program's own \ic{kernel.y} or \ic{fs.y}
beside it shadows std and stays its own business.

\paragraph{\ic{irq N def}.} A top-level declaration that exists only in kernel
mode, whose \ty{N} is the 8259 line (0--15), not the IDT vector:

\begin{lstlisting}
import kernel;

irq 4 def serial_rx() -> void {      // 4 is the 8259 line
    let c = kernel.getc();
    kernel.putc(c);
};
\end{lstlisting}

\noindent The parser records the declaration; the checker requires no
parameters and a \ic{void} result (the generated stub is entered from the
interrupt with nothing to pass and nowhere to put an answer) and refuses a
handler nested inside another \ic{def}; the emitter makes the handler a
body-shake root, because the stub that calls it is generated text the shaker
cannot see; and codegen writes, per declared line, a stub that saves the
caller-saved registers, aligns the stack, calls the handler, acknowledges the
8259 and \ic{iretq}s, and unmasks the declared line. Two handlers on one line,
or one name on two lines, is an error.

The CPU's exceptions (vectors 0--31) get stubs of their own rather than the
halting default: each reports the vector, the error code and, for a page
fault, CR2. A fault is not resumed from.

\subsection{The image and the filesystem}

Y-OS is built as one kernel image and one filesystem image:

\begin{lstlisting}
$ make os            # os/build/os.img and os/build/os.fs
$ make run-os        # boot it here, the serial console on this terminal
$ make run-os-gui    # the same with a screen, keyboard and mouse
\end{lstlisting}

\noindent \ic{make os} needs \ic{qemu-system-x86\_64}. The filesystem is a
YFS1 image; \ic{os/tools/mkfs.y} builds it, \ic{os/manifest.txt} lists what
goes into it and \ic{os/kernel/boot.y} is the kernel program. qemu hands the
kernel one Multiboot module (the filesystem) and it is mounted where it lies:
the files the shell makes and removes live in RAM and are gone when the
machine stops.

\paragraph{YFS1.} Blocks are a superblock field, so the reader is not
hard-wired to a size. Block 0 is the superblock, then the inode table
(64 bytes per inode), then data. An inode has a mode (0 free, 1 file, 2
directory), a size --- bytes for a file, \emph{entry count} for a directory
--- twelve direct block slots and one single-indirect block. A directory entry
is a length, a type, an inode number and a name padded to 4 bytes. An image
must be zeroed where nothing was written, because a free inode is found by
reading a zero mode.

\ic{boot.y} mounts the filesystem, loads \ic{/bin/init} (or, failing that,
\ic{/sh}) into a fresh address space and starts it; when it halts the machine
or exits, so does the kernel. \ic{init} sets the environment
(\ic{HOME}, \ic{USER}, \ic{PATH}, \ic{TERM}, \ic{XDG\_RUNTIME\_DIR}), makes
\ic{/tmp}, \ic{/run/user/1000} and \ic{/home/user}, starts the audio server
\ic{/bin/pulsed} if there is one, runs \ic{/bin/sh} in \ic{/home/user}, and
stays to reap orphans.

\subsection{yunix: the Linux ABI}

\emph{yunix} is the layer that makes the compiler's runtime work on that
kernel: the programs are ordinary \ic{ycc} programs because the kernel
implements the Linux system-call interface the runtime already speaks.
\ic{os/kernel/*.y} are modules of the kernel program \ic{boot.y}, and the
system call table is \ic{sys.y}: \ic{std/kernel.y}'s dispatcher calls
\ic{y\_os\_sys} first for every number, and an answer of
\ic{kernel.SYS\_FALLBACK} hands the call to the kernel's own built-in handler
--- what worked before yunix --- so a stub keeps today's behaviour until its
owner replaces it.

\begin{center}
\begin{tabular}{@{}>{\raggedright\arraybackslash}p{0.24\linewidth}@{\quad}>{\raggedright\arraybackslash}p{0.68\linewidth}@{}}
\toprule
module & what it owns \\
\midrule
\ic{fdt} & per-process descriptor tables: \ic{read}, \ic{write}, \ic{close},
  \ic{lseek}, \ic{ioctl}, \ic{pread}/\ic{pwrite}, \ic{readv}/\ic{writev},
  \ic{poll}, \ic{select}, \ic{dup}/\ic{dup2}/\ic{dup3}, \ic{fcntl} \\
\ic{pipe} & \ic{pipe}, \ic{pipe2} \\
\ic{proc} & \ic{fork}, \ic{vfork}, \ic{execve}, \ic{wait4}, signals (they
  terminate: there are no handlers), ids, \ic{uname}, \ic{kill}, rlimits,
  \ic{reboot} (halt) \\
\ic{time} & a 1\,kHz clock, \ic{clock\_gettime}, \ic{clock\_nanosleep},
  \ic{gettimeofday}, \ic{nanosleep}, \ic{sched\_yield} \\
\ic{thread} & \ic{clone} and \ic{clone3}, \ic{futex}, \ic{arch\_prctl},
  \ic{set\_tid\_address} \\
\ic{mm} & region-based \ic{mmap}, \ic{mprotect}, \ic{munmap}, \ic{brk},
  \ic{madvise}, \ic{getrandom} (\ic{mremap} answers \ic{ENOSYS}) \\
\ic{vfs} & directories and paths over \ic{yfs}: \ic{open}, \ic{openat},
  \ic{stat}/\ic{fstat}/\ic{lstat}/\ic{statx}, \ic{mkdir}, \ic{rmdir},
  \ic{unlink}, \ic{rename}, \ic{chdir}/\ic{getcwd}, \ic{readlink},
  \ic{getdents64}, \ic{truncate}, time stamps \\
\ic{tty} & termios, canonical and raw modes, \ic{VMIN}/\ic{VTIME}, \ic{ISIG},
  window size and job control, over the console \\
\ic{unix} & AF\_UNIX sockets: \ic{socket}, \ic{bind}, \ic{listen}, \ic{accept},
  \ic{connect}, \ic{send}/\ic{recv}, \ic{socketpair}, options \\
\ic{pci}, \ic{ac97}, \ic{dsp} & an AC97 sound card behind \ic{/dev/dsp} with
  the OSS ioctls \\
\ic{tty}, \ic{kbd}, \ic{input} & the console and the input devices \\
\ic{kstat} & the counters the resource monitor and \ic{osstat} read \\
\bottomrule
\end{tabular}
\end{center}

\noindent The hooks the modules plug into live in \ic{std/kernel.y}: the
syscall entry \ic{y\_os\_sys}, the lifecycle calls \ic{y\_os\_task\_forked},
\ic{y\_os\_task\_exec}, \ic{y\_os\_task\_exited} and \ic{y\_os\_task\_switch},
and \ic{block\_restart}/\ic{wake\_chan}/\ic{wake\_expired}. A system call that
must wait does not park in the kernel: it rewinds the user instruction pointer
over the \ic{syscall} and runs again when its channel is woken, so a blocking
read is a normal function that returns when there is something to return.

\paragraph{Memory.} A user program is a Linux executable linked at
\ic{0x400000}, so that address is free in a process's address space. The
kernel maps one 2\,MB page for itself (image, stack, IDT, all below 2\,MB) and
reaches every other physical byte through a direct map at
\ic{0xFFFF888000000000}. Every address space is the kernel's low half, the
direct map, and the process's own user half. An anonymous \ic{mmap} only
reserves a region: the page-fault stub for vector 14 maps a zeroed page on
first touch and resumes the instruction, so a fork copies the few pages a
program touched rather than the whole heap.

\paragraph{The shell and the programs.} \ic{sh} reads a line itself on a
terminal (the tty is raw for the length of the line): cursor and history
editing, Tab completion, quotes and \ic{\bs}, \ic{\$NAME} and \ic{\$?}
expansion, \ic{|}, \ic{<}, \ic{>}, \ic{>>}, \ic{2>}, \ic{;}, \ic{\&\&},
\ic{||}, \ic{\&} and \ic{\#}. Built in are \ic{cd pwd export unset exit env
help true false todesktop} and \ic{NAME=value}; anything else is looked up in
\ic{\$PATH}. Every job runs in its own process group and owns the console
while it runs, so Ctrl-C stops the job and not the shell. Programs in
\ic{/bin} include \ic{ls}, \ic{cat}, \ic{echo}, \ic{mkdir}, \ic{rmdir},
\ic{rm}, \ic{cp}, \ic{mv}, \ic{touch}, \ic{wc}, \ic{head}, \ic{tail},
\ic{pwd}, \ic{stat}, \ic{env}, \ic{clear}, \ic{sleep}, \ic{uname},
\ic{halt}, \ic{write}, \ic{osbench} and \ic{osstat}; a program that reads a
file reads standard input when given none.

\subsection{The screen}

With a display device the machine draws a text terminal and takes keys from a
keyboard and a mouse (\ic{make run-os-gui}, \ic{OS\_VGA=virtio} for
virtio-gpu and virtio-input, or \ic{OS\_VGA=std} for standard VGA with a PS/2
keyboard and mouse). The serial console keeps working beside it, and the
backends follow the host window: a resize gives the terminal new rows and
columns.

\ic{vt.y} is an xterm-style terminal in the kernel (UTF-8 in an 8x16 font,
the 16 colours, a 256-colour cube and 24-bit colour, cursor movement, scroll
regions, the alternate screen, DEC line drawing, mouse reports). The desktop
is what the machine shows at boot: \ic{F1} or \ic{F12} switches between the
terminal and the GUI and the terminal keeps its state while the GUI is up.
\ic{/bin/desktop} shows a wallpaper and a dock whose icons are one per program in
\ic{/bin}; a click types its name into the shell and switches to the terminal.
A file explorer and a resource monitor are panels of the desktop itself. The
wallpaper is a vector image (\ic{os/assets/wallpaper.yvg}) filled by a small
renderer once per size, never per frame. The display is reached through
\ic{/dev/fb0}, which is ioctl-only: information, blit, switch view, type text
and a statistics snapshot.

\subsection{Building, running and testing}
\label{sec:os-cost}

The kernel and system suites are in \ic{make test}:

\begin{itemize}[nosep]
\item the \ic{kernel} suite (\ic{test/kernel/}) builds each fixture with both
  compilers, compares the images, boots them under qemu and checks the serial
  output and the exit status;
\item the \ic{os} suite builds Y-OS with both compilers and requires the
  images to be identical, then boots scripted sessions
  (\ic{test/os/*.cmds}) under qemu and compares the whole serial transcript
  with \ic{test/os/*.expected}. The display has sessions of its own
  (\ic{test/os/display/*.steps}), where each screendump is compared by
  SHA-256 and a mismatch leaves the pictures in \ic{\_check/os-display/}.
\end{itemize}

\noindent \ic{ytest os-session CMD...} scripts a session by hand, and
\ic{YBLESS=1 ./ytest os} rewrites the display hashes after a change that is
meant to move pixels.

\subsection{Limits}

Y-OS is a proof of concept and says where it stops. The kernel heap is a bump
allocator that never gives memory back (64\,MB on a machine with 256\,MB or
more, about 600k path system calls; below that, 8\,MB). Only the 8259
interrupt path is wired: no APIC, no IOAPIC, no MSI. There are no signal
handlers (a signal terminates the task) and no \ic{mremap}. The filesystem and
the terminal live \emph{in} the kernel rather than in user-space servers. The
physical memory allocator hands pages back, but the kernel's own bookkeeping
does not. The GUI has no windows, the screen has no text selection or clipboard
and no font beyond the 8x16 one, and there is no disk: the filesystem is the
boot module. The shell has no job control beyond one foreground job per line
--- no \ic{fg}/\ic{bg}/\ic{jobs}, no Ctrl-Z --- and no shell functions,
\ic{if}/\ic{for}, globbing, history search or wrapped-line editing.
\ic{ycc} itself does not run on Y-OS yet.
Finally, \ic{-g} is refused for a kernel image, because the debug writers are
ELF64/DWARF64.

\paragraph{Cost to a normal compile.} Nothing observable: it is all
behind \ic{target\_is\_kernel()}, the two engines stay byte-identical on the
whole hosted corpus, and the added encoder arms sit at the end of the
instruction ladder, so a hosted line pays for them only if it reaches the end
without matching.

\section{The web target}
\label{sec:web}

\textsf{Y} compiles a second kind of executable beside the hosted ELF and
the bootable kernel: \ic{--target=web} writes a WebAssembly module for
\ic{wasm32}, and the same program runs under Node or in a page.  The target is
not an interpreter and not a bytecode: the compiler's own IR is lowered to
wasm, the context and the token encoding are fixed, and a JavaScript host
supplies the kernel calls the program makes.  ABI 1 names the module's
exports, its metadata section and the host calls; \ic{docs/TARGETS.md} is the
reference table.

\begin{lstlisting}
$ ./ycc --target=web prog.y -o prog.wasm
$ node --disable-wasm-trap-handler src/web/yhost.js run prog.wasm
$ ./ycc --target=web --web-page -o site prog.y   # index.html + prog.wasm
$ ./ycc build --target=web                        # y.toml: build/web/NAME/
\end{lstlisting}

\paragraph{The module.} Values are 64-bit; memory is wasm32 and grows on
demand.  Each Y function body is a \emph{chunk}: a loop with a \ic{br\_table}
on a token and \ic{return\_call}/\ic{return\_call\_indirect} between chunks,
so a call does not grow the wasm operand stack and 1000-deep recursion costs
no host stack.  A token is \ic{2\string^40 + (chunk << 20) + case}, a
64-bit word whose low bits are the case and whose high bits say which
function owns the continuation.  The module exports \ic{memory},
\ic{resume(i64)}, \ic{ctx\_save(i32)}, \ic{ctx\_load(i32)},
\ic{set\_fuel(i32)}, \ic{why()} and \ic{set\_why(i32)}; \ic{resume}
runs until it returns a token, and the host turns that into the next
\ic{resume}, a syscall or a trap.  The custom section \ic{yweb} carries
\ic{abi}, \ic{stack\_lo}, \ic{stack\_hi}, \ic{data\_lo}, \ic{data\_end},
\ic{heap\_lo}, \ic{entry}, \ic{chunks}, \ic{threads}, \ic{gui} and
\ic{audio}.  The initial stack is the Linux one (argc, argv, envp, auxv), so
\ic{argv()} and the environment work unchanged.

\paragraph{The host.} \ic{src/web/core.js} is platform-neutral: it loads and
validates a module, owns the linear memory, the green-thread scheduler, the
syscall table and the host-call table.  \ic{vfs.js} is one VFS for both hosts
(descriptors, pipes, eventfd/timerfd, \ic{/proc}, \ic{/dev}); the Node backend
is the real filesystem, the browser backend an in-memory tree with
\ic{/app} from \ic{files.bin} and \ic{/home/web} in IndexedDB.  Threads are
green threads with a fuel budget: a slice ends after its fuel, a blocking
call parks the thread and the scheduler resumes the next one, and
\ic{pause\_hidden} freezes a hidden page's clock and scheduler.  A raw Linux
syscall is answered by the VFS when it has a plausible answer and refused
with the errno the program would get from a kernel that does not have the
feature; \ic{fork}, \ic{exec}, sockets and \ic{obfuscate} are refused by
design, each with its reason in \ic{test/y/native/*.skip.web}.

\paragraph{The host API.} \ic{y\_os\_host\_sym(name, len)} returns a
positive index for a host function and 0 when there is none;
\ic{y\_os\_host\_call(h, args, n)} calls it with \ic{n} i64 arguments, and a
Promise blocks the calling green thread until it settles.  Strings come back
through a \ic{(buf, cap)} pair: the function returns the length, or
\ic{-(needed)} when the buffer is too small.  The built-in names are
\ic{gui.*}, \ic{audio.*} and \ic{web.*}; a page may register its own under
\ic{app.*}.  \ic{std/web} wraps the page's own API: \ic{available},
\ic{platform}, \ic{sym}, \ic{call}, \ic{call\_str}, \ic{eval}, \ic{log},
\ic{param}, \ic{open\_url} and \ic{fetch}; off the web every one answers
\ic{false}, \ic{""}, \ic{0} or \ic{none}, so the same source still builds and
runs on Linux and Windows.

\paragraph{The screen and sound.} \ic{std/gui} keeps one API.  For web it is
\ic{gui.web.y}, a target overlay merged over \ic{gui.y} and packed as
\ic{gui@web}: the nine \ic{gui.*} calls go through the host symbols, while
present/pump/wait/resize/read\_back/verify\_present/close keep the X11 and
Win32 arms for native builds.  The page owns a canvas: fixed or fit sizing,
integer or automatic zoom (default 1, 100\%), a pixel buffer swapped from
B G R A, and the full
key, char, pointer, wheel, focus, composition and paste input map, with
\ic{ResizeObserver} and \ic{pagehide} becoming events.  Sound is the four
\ic{y\_os\_audio\_*} calls: the browser plays through an inline
\ic{AudioWorklet} (a \ic{ScriptProcessorNode} fallback when a CSP refuses the
Blob URL) and parks a thread while more than 250 ms is queued; Node writes a
WAV with \ic{--audio=FILE.wav}.  \ic{web.available()} is the feature test a
program uses before it asks for a window or a speaker.

\paragraph{Page, site and deployment.} The page provides the stage the
program runs in: one full-window area.  A program that opens a gui window
draws into the viewport, and one without a window gets the console (the page's
stdin/stdout/stderr); \ic{[web] console} = \ic{"show"} or \ic{"hide"} chooses.
A \ic{[web]} table in \ic{y.toml} sets the
title, id, start mode, canvas mode, zoom (100\% unless the app asks otherwise),
console, persistence, \ic{files.bin}, arguments, environment, stack,
pause\_hidden and the single-file form; a \ic{template} points the generator
at the application's own page, whose four placeholders are \ic{TITLE},
\ic{CONFIG}, \ic{HOST} and \ic{MODULE}.  See \ic{docs/TARGETS.md}.

\ic{--web-page DIR} writes \ic{index.html} (the stylesheet and the linked host
are inline) and \ic{NAME.wasm}, plus \ic{files.bin} for \ic{[web] files};
the deployment is those files.  \ic{--web-page=html} additionally inlines the
module and the bundle, so it is one file.  \ic{ycc build --target=web} writes
the same site under \ic{build/web/NAME/}, a single file under
\ic{build/web/NAME.html}, and a Node launcher under \ic{build/web/bin/NAME}
with its private host copy.  A module is fetched as a sibling or as a data URL.
\ic{?ytest=1} writes \ic{YDONE} into the title and \ic{<pre id="ytest">},
which is how the browser suite drives a real page; \ic{?arg=} passes a program
argument and \ic{?fuel=} a slice budget.  The page is inline JavaScript, so a
deployment with a strict CSP must allow it (the module is still fetched, not
eval'd).

\paragraph{What runs.} The whole \ic{test/y/native} corpus runs through the
Node host, and the fixtures that need a system the web does not have skip
with a reason: \ic{fork}, sockets, \ic{exec}, and obfuscation.  Threads,
sleep order, preemption, the collector, NaN and division flags, the token
burst, big bodies, traps, checked builds, stdin, the page API, gui and sound
all have web fixtures.
Performance and size are measured, not guessed: the wasm module is several
times a native build's work but a browser runs it inside a frame, and the
network cost is the module's gzip size.

\paragraph{Limits.} A web build has no \ic{fork}, \ic{exec}, \ic{popen},
sockets, \ic{/proc/self/mem} or \ic{ptrace}; a program that needs a real
kernel stays native.  \ic{--pack}, \ic{--obfuscate}, \ic{--target=kernel} and
\ic{-g} either refuse the target or are not part of it.  The host is not a
security boundary: a page's JavaScript can read its own module and memory as
any page can, and a save in IndexedDB is the player's.

